\section{Вывод}

В ходе выполнения лабораторной работы были успешно реализованы:

\begin{enumerate}
    \item \textbf{Динамические библиотеки}, реализующие заданный функционал в двух вариантах
    \item \textbf{Программа №1}, использующая статическое связывание с библиотекой на этапе компиляции
    \item \textbf{Программа №2}, осуществляющая динамическую загрузку библиотек во время выполнения
\end{enumerate}

\textbf{Основные выводы:}

\begin{itemize}
    \item Динамические библиотеки позволяют эффективно организовывать код, разделяя его на переиспользуемые компоненты
    \item Статическое связывание обеспечивает простоту распространения программ, но увеличивает их размер
    \item Динамическая загрузка предоставляет гибкость в управлении функционалом программы без её перекомпиляции
    \item Оба подхода имеют свои области применения: статическое связывание предпочтительно для простых приложений, динамическая загрузка — для модульных и расширяемых систем
\end{itemize}

Полученные практические навыки работы с динамическими библиотеками будут полезны при разработке сложных программных систем, требующих модульности и гибкости архитектуры.

